\section{Preparation, kernels, and projective holes}
\label{prep:section}

This section supplies the two kinds of control needed for projective
windows. Preparation makes the boundary limits finite and preserves
their real parameters. The second half of the section proves that
projective fibres which are present up to algebraic holes retain this
property throughout a connected parameter domain.

All manifolds in this section are complex manifolds unless real coordinates
are specified. A locally pseudoconvex open subset of a manifold means an
open subset which is locally Stein at every boundary point. The open sets are
allowed to be disconnected before a distinguished component is selected.
Subanalytic data on a relatively compact coordinate box are always assumed
to extend as globally subanalytic data to a slightly larger box. Rational
powers of a positive scale variable are allowed. Polynomial descriptions
in projective variables are understood in a finite collection of bounded
algebraic charts, including charts at infinity. In this paper,
\emph{bounded algebraic data} means that one such finite atlas and
uniform finite bounds for the number and degrees of the polynomial
fibre equations and tests have been fixed. The coefficient arrays
depend analytically, or globally subanalytically, on the retained
parameter boxes as specified in each statement. For an algebraic
map, the same convention is imposed on its graph. A \emph{bounded-data
semialgebraic family} has a fixed finite Boolean description of this
kind in real and imaginary fibre coordinates. These are bounds on
finite algebraic complexity; estimates for coefficients, scale
matrices or their inverses are stated separately when required.

\subsection{The preparation conventions}

\begin{definition}[Generic real parameters]\label{prep:generic}
An assertion at generic parameters in an open real box means that it holds
outside a locally finite union of subanalytic subsets of smaller dimension.
When countably many fixed compact boxes or positive accuracies are tested,
a countable union of such exceptional sets is permitted. Each application
involving finitely many data is subsequently restricted to a neighborhood
of the chosen parameter on which those data have the asserted regularity.
No uniform neighborhood for all members of a countable list is implicit.
\end{definition}

We use the following standard forms of real analytic preparation.
The underlying results are the uniformization and rectilinearization
theorems and the distance inequality of
\cite[Theorems 0.1, 0.2, and 6.4]{BierstoneMilman1988}, together with
one-variable preparation with parameters
\cite[Theorem 1 and Sections 0.3--0.4]{LionRolin1997}.
Parusi\'nski's earlier preparation theorem
\cite[Theorem~7.5]{Parusinski1994} is a predecessor of this method;
the specified centre condition used below is the one in Lion--Rolin.
For finite definable partitions, choice and components we use
\cite[Theorems~2.10, 3.1 and~3.9]{Coste1999}.

\begin{lemma}[Power bounds and generic Puiseux preparation]
\label{prep:puiseux}
Let $A$ be globally subanalytic and let $f(v,t)$ be globally subanalytic
for $v\in A$ and $0<t<\epsilon(v)$, where $\epsilon$ is positive.
If $f(v,t)>0$, there is an integer $N$ such that, for every fixed $v$,
\[
 t^N<f(v,t)<t^{-N}
 \quad\hbox{for all sufficiently small }t>0.
\]
The threshold may depend on $v$. For finitely many such functions one
may use a common $N$.

If $A$ is an open real box, then near a generic $v_0$ a given finite
collection of functions admits convergent expansions
\[
 f(v,t)=\sum_{j\geq j_0}a_j(v)t^{j/q}.
\]
with a common positive integer $q$ and coefficients real analytic in $v$.
After multiplication by a sufficiently high power of $t$, these are
jointly real analytic functions of $(v,t^{1/q})$ through $t=0$.
Convergence and every fixed number of $v$-derivatives are uniform on
smaller compact parameter boxes.
\end{lemma}

\begin{proof}
Apply one-variable preparation with $t$ as the last variable. The finitely
many preparation cylinders meeting $t=0$ have, on their bottom bands,
center zero: a nonzero center comparable to $t$ cannot remain comparable
as $t$ tends to zero with $v$ fixed. On each such band the prepared
function is a rational power of $t$ times a nonzero function of $v$
times a unit bounded above and below. There are only finitely many
rational exponents. Taking $N$ larger than all their absolute values
absorbs the positive factors depending on fixed $v$ and proves the
first assertion. Apply the same argument to $1/f$ when necessary.

For the second assertion, restrict to a full-dimensional base cell on
which the finitely many coefficient functions in the preparation are
analytic and their required nonzero denominators remain nonzero.
In the analytic unit choose projective coordinate charts about the
limiting arguments. An argument involving a negative power of $t$
either has identically zero numerator on the chosen base cell or
tends to infinity; in the latter case its reciprocal is a valid
analytic coordinate and its nonzero coefficient can be inverted.
In these charts the arguments are analytic in $v$ and nonnegative
rational powers of $t$.
Clear their denominators. The unit is analytic on a neighborhood of
its compact argument set, giving the asserted joint analytic extension.
Ordinary convergence of analytic power series on a smaller box gives
the derivative assertion. The discarded base cells and coefficient
singularities have smaller dimension.
\end{proof}

\begin{remark}\label{prep:fixed-parameters}
An estimate obtained with additional parameters held fixed does not
permit those parameters to move during a correction. When a distance
inequality is needed fibrewise, apply preparation to the definable
fibrewise distance and residual functions. This gives a common power
on finitely many parameter pieces; its positive constants and thresholds
may still depend on the fixed parameters. Joint analytic preparation,
when needed, requires restriction to a full-dimensional parameter piece.

Here is an explicit way to retain the fixed parameter. For compact
closed fibres $C_v$ and a nonnegative continuous definable residual
$r(v,z)$ whose zero set is $C_v$, put
\[
 m(v,s)=\min\{r(v,z):\operatorname{dist}(z,C_v)\geq s\}
\]
on a compact test box, assigning the value $1$ when the set is empty.
This is positive for $s>0$. Lemma \ref{prep:puiseux} bounds it below
by $s^N$ for sufficiently small $s$, with $N$ independent of fixed
$v$. Compactness excludes points of fixed positive distance when
their residual tends to zero. Consequently, for small residuals,
\[
 \operatorname{dist}(z,C_v)\leq r(v,z)^{1/N}.
\]
All corrections therefore take place in the same fibre $C_v$.
Empty $C_v$ cause no spurious limit, since the residual then has
a positive minimum on the compact test box.
\end{remark}

\begin{lemma}[Projective models with parameters]\label{prep:projective-models}
Suppose an algebraic family in projective coordinates is pulled back
by a local analytic parameter map. Suppose an irreducible member is
specified by a marked irreducible germ. Near a generic point of a
maximally totally real open parameter box, one may track that component
on a local branch and replace the family and its evaluation maps by
a smooth proper projective family of connected smooth models.
The evaluation maps are holomorphic, and their restrictions to the
projective fibres are algebraic of bounded data. Real analytic
parameter dependence can first be complexified.
\end{lemma}

\begin{proof}
Track irreducible components on a relative Hilbert or Chow parameter
space; for their algebraic existence and projectivity see
\cite[\S3, Theorems~3.1--3.2]{Grothendieck1961} and
\cite[\S2.4(b), Proposition~4]{AndreottiNorguet1967}. On the locus of a fixed component type the tracking map is
generically finite. Remove its branch and nonflat loci and choose a
local branch. Resolve the resulting total family and the graphs of
the finitely many evaluation maps by projective modifications. Generic
smoothness removes a further proper analytic subset of the complex
parameter space. A maximally totally real open box is not contained
in a proper complex analytic subset, by the holomorphic identity
theorem. Thus these deletions are legitimate at generic real parameters.
The selected generic fibres are irreducible, and their smooth models
are connected. After shrinking, properness and smoothness preserve
connectedness. All constructions are made with finitely many projective
equations and modifications; restricting to a relatively compact
parameter box supplies the stated bounded data.
\end{proof}

\subsection{Open kernels and finite jets}

\begin{definition}[Open kernel]\label{prep:kernel}
Let $U_t$ be open subsets of a fixed manifold $M$, $0<t<t_0$. Their
open kernel is
\[
 \mathcal K(U_t)=\{z\in M: \text{some neighborhood }W\ni z
       \text{ satisfies }W\subset U_t
       \text{ for every sufficiently small }t\}.
\]
Equivalently, if $F_t=M\setminus U_t$, then
\[
 \mathcal K(U_t)=M\setminus\limsup_{t\downarrow0}F_t,
\]
where the upper limit allows both $t$ and the point of $F_t$ to vary.
In a varying family these tests are relative to the total space;
smooth local trivializations identify the limiting fibre. A marked
kernel is the connected component containing a specified point of
the open kernel.
\end{definition}

The neighborhood quantifier in this definition is essential. In
particular, pointwise eventual membership is insufficient. The compact
subsets of an open kernel are eventually contained in $U_t$, with a
neighborhood, by a finite covering argument.

We use the Hartogs-figure criterion for local pseudoconvexity and the
Euclidean Levi theorem in their standard forms; see
\cite[II.3.4 and II.3.7]{FritzscheGrauert2002} and
\cite[VIII.9.11(a)]{Demailly2012}. The exhaustion and directional-radius
criteria used below are \cite[I.7.2]{Demailly2012}.

\begin{lemma}[Pseudoconvexity of open kernels]
\label{prep:kernel-pseudoconvexity}
An open kernel of locally pseudoconvex open sets is locally
pseudoconvex. The assertion also holds in varying manifolds equipped
with smoothly converging holomorphic coordinate charts.
\end{lemma}

\begin{proof}
Use the local Hartogs-figure criterion in a coordinate polydisc. A
relatively compact Hartogs figure with a neighborhood contained in
the kernel is contained, together with that neighborhood, in every
$U_t$ for sufficiently small $t$. Pseudoconvexity fills its associated
polydisc in each $U_t$. To obtain neighborhood inclusion for a compact
subset of the filling, slightly enlarge the figure within the given
neighborhood before applying the criterion. The resulting enlarged
fillings supply a neighborhood of that compact subset in every small
$U_t$. Thus the filling lies in the kernel. This is precisely the
Hartogs-figure criterion for the kernel. The same argument uses the
corresponding figures in smoothly varying holomorphic charts; strict
compact containment persists under this change.
\end{proof}

\begin{lemma}[Finite-jet specialization]\label{prep:finite-jet-specialization}
In bounded projective charts let a family of sets be described by
finitely many Boolean combinations of equations and inequalities
polynomial in the real fibre coordinates $y$, with coefficients real
analytic in real parameter coordinates $b$. Fix $x$ and substitute
\[
 b=x+R_t h,
\]
where the entries of $R_t$ have convergent Puiseux expansions,
$R_t\to0$, and both $R_t$ and $R_t^{-1}$ have power bounds.
The upper limits of these sets, and hence the open kernels of their
complements, are semialgebraic in $(h,y)$ on the central fibre.
Here $h$ ranges over its whole affine space, while the substitution
is tested only where $b$ belongs to the original parameter box.
When external parameters are present, assume that all data, including
$R_t$, are jointly globally subanalytic and that the displayed
analytic coefficients are jointly analytic. Then the resulting
descriptions have analytic coefficients after generic restriction.
If all the original data are merely
globally subanalytic, the corresponding conclusions are subanalytic.
\end{lemma}

\begin{proof}
The subanalytic assertion follows immediately by writing the
upper-limit quantifiers with a positive accuracy and a positive
upper bound for $t$. We prove the stronger assertion in the stated
polynomial-in-$y$ situation. Work on one basic piece, written as
\[
 F_i(b,y)=0,\qquad H_j(b,y)\geq0,\qquad L_k(b,y)>0.
\]
Finite unions cause no difficulty. First require a common strict
slack $L_k\geq t^A$. This does not change the upper limit when $A$
is sufficiently large. Indeed, for a proposed limit point, a positive
accuracy, and a fixed bound on $h$, take the supremum of the minimum
strict slack among witnesses within that accuracy; cap this supremum
by $1$. Whenever witnesses exist for arbitrarily small $t$, they
exist for every sufficiently small $t$ by one-dimensional
subanalyticity. The positive supremum is definable. Lemma
\ref{prep:puiseux} gives one exponent $A$ for this family, allowing
the threshold to depend on the point, accuracy, and bound. A witness
with at least half the supremum gives the desired slack after
increasing $A$.

For a second exponent $B>A$, allow residual errors of size $t^B$
in the equations and weak tests. Include the slack as an extra
variable. On a slightly larger compact original-coordinate box the
distance inequality for the resulting closed semianalytic set gives
\[
 \operatorname{dist}((b,y,\sigma),C)
      \leq c\,\operatorname{residual}(b,y,\sigma)^{\theta}
\]
for some $\theta>0$. The same statement with external parameters
fixed follows from the fibrewise version described in Remark
\ref{prep:fixed-parameters}. Choose $B$ so large that
$B\theta>A+1$ and
$\|R_t^{-1}\|t^{B\theta}\to0$.
A point satisfying the relaxed tests can then be corrected to $C$;
the correction is $o(t^A)$ in slack, $o(1)$ in $y$, and $o(1)$ in
the normalized coordinate $h$. Strict tests still have, for example,
half the original slack. Thus these relaxed tests have exactly the
original upper limits.

Taylor-expand every coefficient at $b=x$. Since $R_t\to0$ with a
positive power bound, sufficiently high finite Taylor order makes
all remainders $o(t^B)$ on each fixed bounded $h$-box. Truncate the
Puiseux series at correspondingly high order. Changing the relaxed
error by a fixed factor absorbs these remainders and leaves the same
upper limit, by the preceding correction argument. The resulting
tests are polynomial in $h,y$ and in a common root of $t$, with
finitely many negative powers removable by multiplication by a
positive power of $t$. Real quantifier elimination therefore gives
a semialgebraic upper limit.

The exponent choices can be made independently of the fixed bound
on $h$: include its reciprocal as another bounded fixed parameter
in the slack and distance tests. The constants and the small-$t$
threshold may deteriorate with that bound; the finitely many power
exponents do not. Consequently one finite polynomial description
works on the full affine $h$-space. A finite projective atlas treats
all fibre points. The finite Taylor coefficients and the finitely
many choices used above are analytic in external parameters after
generic restriction. In particular, apply Lemma \ref{prep:puiseux}
jointly to the entries of $R_t$ and $R_t^{-1}$ and discard the zero
loci of their nonzero leading coefficients. Their relevant positive
orders of decay and inverse power bounds are then fixed on the
smaller parameter box. This supplies a common Taylor order there,
rather than only a separate order for each fixed parameter.
\end{proof}

\begin{lemma}[Generic convergence of distance tests]
\label{prep:generic-distance-convergence}
Let $x$ range over an open real box $U$, let $K$ be compact and globally
subanalytic, and let
$d_t(x,z)$ be a bounded globally subanalytic family of functions on
$K$, for $0<t<t_0$, uniformly Lipschitz in $z$. There is a relatively closed
subanalytic subset $E\subset U$ of dimension smaller than $\dim U$
such that the extension
\[
 \widehat d(x,z,t)=
 \begin{cases}
 d_t(x,z),&t>0,\\
 d_0(x,z),&t=0,
 \end{cases}
 \qquad d_0(x,z)=\lim_{t\downarrow0}d_t(x,z)
\]
is jointly continuous at every $(x,z,0)$ with $x\in U\setminus E$
and $z\in K$. In particular, for such $x$, if $x_t\to x$ and
$z_t\to z$, then
\[
 d_t(x_t,z_t)\longrightarrow d_0(x,z).
\]
For every compact $L\subset U\setminus E$, the convergence
$d_t\to d_0$ is uniform on $L\times K$. One common exceptional
set and one smaller open center box may be used for finitely many
such tests. For countably many tests one may exclude the union of
their exceptional sets, but no common open center box is asserted.
These conclusions apply in particular to bounded, capped distance
functions of definable closed sets.
\end{lemma}

\begin{proof}
Pointwise limits exist by one-dimensional definable monotonicity.
They have the same Lipschitz bound in $z$. A finite net in $K$ thus
upgrades convergence at fixed $x$ to uniform convergence on $K$.
For each positive rational $\delta$, the supremum of the positive
numbers $a$ for which
\[
 |d_t(x,z)-d_0(x,z)|<\delta
       \quad(0<t<a, z\in K)
\]
is a positive definable function of $x$, after capping it by $1$.
It is continuous on a dense open definable set, and hence locally
bounded below there. This controls the error also at moving centers
$x_t$ near a chosen generic $x$. For each point of a countable dense
subset of $K$, the definable function $d_0(x,z)$ is likewise
continuous at generic $x$. A finite-net argument and the common
Lipschitz constant give continuity in $x$ uniformly on $K$ at all
points outside the resulting countable union of lower-dimensional
subanalytic sets. Combining the two estimates proves that
$\widehat d$ is continuous at every point of $\{x\}\times K\times\{0\}$
for all such $x$.

To obtain one open good locus for this fixed compact test, define
\[
 \Sigma=\{x\in U:\text{ for some }z\in K,\
        \widehat d\text{ is not continuous at }(x,z,0)\}.
\]
This is a subanalytic set: both $d_0$ and the extension are
definable, and failure of continuity is expressed by the finite
quantifiers
\[
 \exists z\in K\ \exists\epsilon>0\ \forall\delta>0\
 \exists(x',z',t'):\quad
 \begin{cases}
 x'\in U,\ z'\in K,\ t'\geq0,\\
 |x'-x|+|z'-z|+t'<\delta,\\
 |\widehat d(x',z',t')-d_0(x,z)|\geq\epsilon.
 \end{cases}
\]
All variables are restricted to the domains of the given functions.
The preceding argument puts $\Sigma$ inside a countable union of
lower-dimensional subanalytic sets. Their closures are nowhere
dense, so the Baire theorem shows that $\Sigma$ has no interior.
Since $\Sigma$ itself is subanalytic, it has dimension smaller
than $\dim U$, as does its relative closure $E$.

On $U\setminus E$ the extension is therefore jointly continuous
at every boundary point with $t=0$. Compactness of $L\times K$
turns this into uniform convergence on that product. Taking a
finite union of the exceptional sets proves the finite-family
assertion. Taking a countable union gives precisely the weaker
countable-family assertion stated above.
\end{proof}

\begin{corollary}[Freezing the real center]\label{prep:generic-freezing}
Let a definable family be written using $\operatorname{Re}b$ and
$\operatorname{Im}b$, and let $R_t\to0$ be a fixed real power matrix,
independent of the real center $x$. For the substitution
\[
 b=x+R_t(u+iv)
\]
its upper-limit and open-kernel tests at generic $x$ are independent
of $u$. Thus the kernels are invariant under real translations in
the scaled coordinates. The tests may include remaining bounded
fibre variables jointly. More generally, when a normalization
$R_t(x)$ depends definably on $x$, Lemma
\ref{prep:generic-distance-convergence} applies to the normalized
closed-set tests themselves.
\end{corollary}

\begin{proof}
For the fixed matrix use $a=\operatorname{Re}b$ as an independent
parameter and take the closed-set distance tests after setting
$\operatorname{Im}b=R_tv$. Their generic limit at $a=x$ is unchanged
by replacing $a$ by $x+R_tu$, by Lemma
\ref{prep:generic-distance-convergence}. This applies to arbitrary
bounded $u,v$ and to their convergent perturbations, which are the
quantifiers defining upper limits. Taking complements gives the
kernel assertion. The final statement is the same lemma applied
after, rather than before, the parameter-dependent normalization.
\end{proof}

\begin{lemma}[Power diagonalization]\label{prep:power-diagonalization}
Let $f(e,t,v)$ be a bounded globally subanalytic function with
iterated limit $f_{00}(v)=\lim_{t\downarrow0}\lim_{e\downarrow0}
f(e,t,v)$. There is an integer $N_0$ such that for every integer
$N>N_0$,
\[
 \lim_{t\downarrow0}f(t^N,t,v)=f_{00}(v)
\]
for each fixed $v$. The same choice works for finitely many tests;
the parameter $v$ may include fixed box sizes and positive
accuracies. For bounded closed-set tests, apply this to the distance
functions on slightly larger boxes, retaining a smaller box to
avoid artificial edge effects.
\end{lemma}

\begin{proof}
Let $f_0(t,v)=\lim_{e\downarrow0}f(e,t,v)$. For fixed positive
$\delta$, the supremum of those $a\in(0,1)$ such that
$|f(e,t,v)-f_0(t,v)|<\delta$ for every $0<e<a$ is positive and
definable. Lemma \ref{prep:puiseux}, with $(v,\delta)$ held fixed,
bounds this threshold below by $t^{N_0}$ for all sufficiently small
$t$, with $N_0$ independent of $(v,\delta)$. The asserted diagonal
then lies below the threshold. Let $t$ tend to zero and then let
$\delta$ tend to zero.
\end{proof}

\subsection{Polynomial tests for compact holes}

We now turn from limits of parameter sets to omissions inside compact
fibres. The basic test measures a polynomial on the omitted set.
Its plurisubharmonicity will turn containment along a parameter edge
into containment throughout a neighbourhood.

\begin{lemma}[Polynomial maximum over the holes]
\label{holes:polynomial-maximum}
Let $B$ be a complex manifold and let $J\subset B\times\mathbb C$
be closed. Assume that $J\to B$ is proper, its fibres are locally
bounded in the line coordinate, and $(B\times\mathbb C)\setminus J$
is locally pseudoconvex. If $g(b,t)$ is monic of positive degree in
$t$ with holomorphic coefficients in $b$, then
\[
 \psi_g(b)=\log\max_{t\in J_b}|g(b,t)|
\]
is plurisubharmonic, with value $-\infty$ for an empty fibre or a
zero maximum. The assertion also holds for polynomials with
nowhere-zero holomorphic leading coefficient, and for the constant
test $1$.
\end{lemma}

\begin{proof}
Work over a small Stein coordinate box. The map
$F(b,t)=(b,g(b,t))$ is finite, proper, and open. Put $J'=F(J)$.
Its complement $\Omega'$ is locally pseudoconvex. Here is the
justification, including the ramification locus. Take a
plurisubharmonic exhaustion of
$\Omega=(B\times\mathbb C)\setminus J$, which is Stein over the
small box by the local Levi criterion. On $\Omega'$ take the maximum
of this exhaustion over all preimages under $F$. Away from branch
values this is a finite maximum of plurisubharmonic functions.
Local boundedness above and continuity of the finite fibres extend
it plurisubharmonically across branch values. It is an exhaustion:
a sequence approaching $J'$ has at least one preimage approaching
$J$, while properness of $F$ handles escape to infinity. This proves
the claim about $\Omega'$.

Invert the image line coordinate and adjoin its point at infinity.
In the inverse coordinate $w$, the resulting domain contains
$w=0$ over every base point. It is locally pseudoconvex there by
local boundedness of $J'$. Its largest centered vertical disc has
radius
\[
 r(b)=\bigl(\max_{v\in J'_b}|v|\bigr)^{-1},
\]
where $r=+\infty$ when the maximum is zero or the fibre is empty.
The Hartogs radius criterion says that $-\log r$ is
plurisubharmonic, proving the assertion. Multiplication by a
nowhere-zero holomorphic leading coefficient adds a pluriharmonic
function.

For the constant test, note first that empty fibres form an open
set by properness. Choose, locally on the base, a constant $a$ outside
a disc containing every $J_b$. The already proved test $t-a$ has
a positive lower bound on every nonempty fibre and is $-\infty$
on empty fibres. If there is an empty fibre, it is $-\infty$ on an
open base set and therefore identically $-\infty$ on the connected
base box. Thus either every fibre is empty or none is. The test
for $1$ is accordingly constant $-\infty$ or constant zero.
\end{proof}

In this proof the Hartogs radius criterion concerns the largest
disc centered at a fixed holomorphic section; the whole vertical
section need not be a disc. Equivalently, take its complete circular
Hartogs subdomain. The criterion follows from the local
Hartogs-figure test, and gives the negative logarithm of that radius.

\begin{lemma}[Boundary trapping]\label{holes:boundary-trapping}
Let a proper holomorphic submersion be given near a relatively compact
parameter interval $I$ on the boundary of a one-sided complex disc.
Let $G$ be locally pseudoconvex on that side, and suppose every
limit of its fibre holes as the transverse parameter tends to $I$
belongs to a relative divisor $D$ that extends holomorphically
across $I$ and is proper in each fibre. Then, after restriction near
an interior point of $I$, the holes on that side belong to $D$.
The conclusion holds with additional fixed parameters whenever
these hypotheses hold locally uniformly in those parameters.
\end{lemma}

\begin{proof}
At a possible limiting hole choose fibre coordinates $(z,t)$ in
which $D$ is given by a Weierstrass polynomial $g(s,z,t)$. Choose
the vertical disc with boundary disjoint from $D$ and shrink all
other variables. The limit hypothesis confines the holes to a
strictly smaller vertical disc. After filling trivially outside
this disc, they give a closed set $J$ as in Lemma
\ref{holes:polynomial-maximum}, with base $(s,z)$.

For fixed $z$, the subharmonic function $\psi_g(s,z)$ is bounded
above near a compact subinterval of $I$ and tends locally uniformly
to $-\infty$ there. To see that it is identically $-\infty$, fix a
smaller half-disc. Harmonic measure comparison bounds it by
$C-M\omega(s)$ for every $M>0$, where $\omega(s)>0$ is the
harmonic measure of a nontrivial subinterval. Let $M$ tend to
infinity. Thus $g$ vanishes at every hole. Finitely many such
cylinders cover the possible holes on compact fibre sets; outside
them the limit hypothesis already excludes holes. This proves
the assertion, including its locally uniform version.
\end{proof}

\begin{lemma}[Interior trapping]\label{holes:interior-trapping}
Let $G$ be open and locally pseudoconvex over a whole connected
base box, with the preceding local confinement hypotheses.
Suppose a relative
divisor $D$ is holomorphic in a whole base box and contains the
holes on a nonpluripolar parameter slice inside the box. Then it
contains the holes nearby. In particular this applies to an open
piece of a real analytic maximally totally real slice.
\end{lemma}

\begin{proof}
Use the same Weierstrass cylinders and polynomial maximum. For
fixed remaining fibre coordinates, $\psi_g$ is $-\infty$ on the
specified nonpluripolar set. A plurisubharmonic function with this
property is identically $-\infty$ on its connected box. For the
last stated case this can also be checked without a capacity
criterion: straighten the real analytic slice biholomorphically,
and apply the one-variable uniqueness principle successively to
the real intervals in a real coordinate box. The resulting
vanishing on the complex box proves containment.
\end{proof}

\begin{lemma}[Tangency propagation in a connected fibre]
\label{holes:tangency-propagation}
Let $q:Q\to\Delta$ be a holomorphic submersion, let $G\subset Q$
be open and locally pseudoconvex, and let $A$ be a divisor in
$Q_0$. Suppose $Q_0\setminus A$ is connected. Assume that on one
side of a regular embedded real analytic arc through $0$, every
limit of holes belongs
to $A$. If $G_0$ is nonempty, then
\[
 Q_0\setminus A\subset G_0.
\]
It suffices to have these hypotheses on neighborhoods of compact
paths in $Q_0\setminus A$.
\end{lemma}

\begin{proof}
Since $G_0$ is open and $A$ has no interior, it contains a point
outside $A$. Connect that point to any prescribed point of
$Q_0\setminus A$ by a path and cover the path by finitely many
submersion boxes avoiding $A$ on the central fibre. On a sufficiently
small fixed one-sided disc the limit hypothesis and compactness
give all these boxes on the good side. In each box choose a disc
$\Delta(c,r)$ in the transverse parameter internally tangent to
the arc at $0$. The full product of this disc with a smaller
connected vertical patch belongs to $G$.

For a vertical point $z$ let $R(z)$ be the largest radius of the
transverse disc centered at $c$ contained in $G$ within this box.
The Hartogs radius criterion gives that $-\log R$ is
plurisubharmonic, and $R\geq r$ everywhere. At the starting patch,
openness of $G$ at the tangency point gives $R>r$ at some point:
the rest of the closed tangent circle is compactly on the good
side. If $R=r$ anywhere in the connected patch, $-\log R$
attains its maximum there and the strong maximum principle
would force $R=r$ everywhere. Thus $R>r$ throughout the patch,
so the central-fibre points of that patch are in $G$. Continue
through successive overlapping boxes along the path. This reaches
the prescribed point.
\end{proof}

\subsection{Thin holes and propagation in a projective family}

The preceding lemmas propagate one specified container. We next show
that bounded algebraic containers can be chosen through a boundary
wall, then use this to propagate the existence of a Zariski open
fibre. Analyticity of a sufficiently thin omitted set completes the
argument across the remaining lower-dimensional parameter strata.

\begin{theorem}[Analyticity of a thin pseudoconcave set]
\label{holes:thin-analyticity}
Let $M$ be a complex manifold of dimension $n$, and let $A\subset M$
be closed and subanalytic, with real dimension at most $2n-2$.
If $M\setminus A$ is locally pseudoconvex, then $A$ is either
empty or an analytic set of pure complex codimension one.
\end{theorem}

\begin{proof}
The measure hypothesis in the standard analyticity theorem is
automatic here. Indeed uniformize the closed subanalytic set near
a compact coordinate box by a proper analytic map from a manifold
of dimension at most $2n-2$. The inverse image of the box is
compact. A finite covering by coordinate patches makes the map
Lipschitz on compact smaller patches, so its image has finite
$(2n-2)$-dimensional Hausdorff measure there. The local Stein
complement means exactly that $A$ is $1$-concave in the sense of
\cite[Definition 1]{Vajaitu2000}. For $n>1$, apply
\cite[Theorem 1]{Vajaitu2000} with $q=1$; it gives analyticity and
pure dimension $n-1$. For $n=1$, a zero-dimensional closed
subanalytic set is locally finite and hence is an analytic divisor
directly. For smooth ambient spaces, the thin-pseudoconcave-set method
already appears in \cite[Theorem~1.4]{Hirschowitz1973}.
\end{proof}

\begin{lemma}[A relative divisor containing a proper subset]
\label{holes:relative-divisor-container}
Let $q:Q\to B$ be a smooth projective family with connected fibres,
and let $A\subset Q$ be a relative algebraic subset with analytic
coefficient data, proper in every fibre. Near a generic parameter
there is a relative effective divisor $D$ containing $A$, defined
over the full smaller parameter box and proper in every fibre.
\end{lemma}

\begin{proof}
First restrict near a generic parameter at which $A\to B$ is flat.
Since $q$ is smooth, the exact sequence of its ideal $\mathcal I_A$
shows that $\mathcal I_A$ is flat over this smaller base as well.
In a relative projective embedding, choose a twist
$\mathcal I_A(k)$ which is generated and has vanishing higher
cohomology on the selected fibre
\cite[Tag~0B5T]{StacksProject}, using the comparison in
\cite[Theorems~1--3]{Serre1956}. Upper semicontinuity preserves
the vanishing after shrinking; constancy of the Euler characteristic
in a proper flat family makes $h^0$ constant. Grauert's
constant-cohomology theorem therefore gives a locally free direct
image and fibrewise base change
\cite[\S7, S\"atze~3, 5 and~6]{Grauert1960}; see the correction
\cite[p.~36]{Grauert1963} for the distinct general base-change statement.
The evaluation map generates on the selected fibre. Nakayama's
lemma and properness allow a further shrinking on which it generates
throughout. We thus have local holomorphic sections of the twisted
ideal whose restrictions generate every fibre ideal. Choose one whose
restriction at the selected parameter is not identically zero;
such a section exists because $A$ is proper in that fibre.
Extend it over a smaller base box. It stays nonzero at a local
section of $q$ through a point where its central restriction is
nonzero. Thus its zero divisor is proper in every nearby fibre
and contains $A$. The construction uses one fixed twist, hence
has bounded algebraic fibre degree. Generic flatness supplies the
initial parameter outside a proper analytic subset; the subsequent
choices only shrink its neighbourhood. Such a parameter can be
chosen in the maximally totally real box. This argument does not
require the general bad base-change locus to be closed.
\end{proof}

\begin{lemma}[Bounded algebraic containers]\label{holes:bounded-containers}
Consider a smooth proper projective family with connected fibres, and
a bounded-data semialgebraic family of subsets of those fibres.
There is a degree bound, depending only on the
family description, such that every member contained in a proper
complex algebraic subset is contained in an effective divisor of
that degree or smaller. The parameters admitting such a container
form a subanalytic set, and containers can be selected subanalytically.
For a one-sided family approaching a generic real analytic parameter
wall, its limiting containers can be chosen as proper effective
divisors varying real analytically on that wall.
\end{lemma}

\begin{proof}
Regard the finitely many real polynomial coefficients and the
projective fibre equations as independent algebraic parameters.
Real algebraic cell decomposition, with these parameters first,
expresses each fibre as finitely many Nash pieces of uniformly
bounded algebraic complexity. Their algebraic closures, and the
complex Zariski closures of their projections into the complex
fibre coordinates, consequently have bounded degree by elimination.
The complex Zariski closure of the original set is the union of
these finitely many closures. When it is proper, a homogeneous
polynomial of bounded degree vanishes on it without vanishing
identically on the ambient fibre. Its zero divisor is a container.

At a fixed degree the condition that a homogeneous equation
vanish on every point of the semialgebraic set, but not on the
whole fibre, is a quantified semialgebraic condition on its
coefficients. Quantifier elimination and definable choice give
the parameter locus and the selection. In a varying smooth
projective family, record the chosen zero divisors as cycles in
the finite union of relative Chow spaces of these degrees.
Their closure is proper over a smaller parameter box. Its
limiting cycles retain their fibre dimension, so their supports
are proper divisors; a limiting homogeneous equation alone would
not suffice, since it could vanish identically on a special fibre.

Across a generic real wall, one-sided subanalytic limits of this
bounded cycle data exist and vary real analytically after a
stratification and generic restriction. This follows equally by
applying Lemma \ref{prep:puiseux} to the normal wall variable in
projective coordinate charts. More precisely, restrict to a
full-dimensional real analytic stratum of the wall on which the
cycle degree and these charts are fixed. Preparation in the
one-sided normal variable gives an analytic limit map on that
stratum and convergence uniform on each smaller compact wall
box. Restrict this map to the real interval of a transverse
analytic complex disc. Its real analytic coefficients complexify
holomorphically in the disc variable, jointly real analytically
in the other wall parameters. The equations defining the relative
Chow space vanish on the real interval, so continue to vanish by
the holomorphic identity theorem. Pulling back its universal cycle
gives a relative divisor on the smooth pulled-back family.
The analytic extensions exist on a common small disc after
restriction to a compact smaller wall box. This is the claimed
locally uniform continued container; no holomorphic extension in
all complex base variables at once has been asserted.
\end{proof}

\begin{proposition}[Projective-fibre propagation]
\label{holes:projective-propagation}
Let $q:Q\to B$ be smooth, proper, and projective, with connected
smooth fibres. Let $G\subset Q$ be open, subanalytic, locally
pseudoconvex, and fibrewise semialgebraic with bounded data.
Put $S=q(G)$, and let $S_0$ be a connected component of $S$.
If $G_b$ contains a Zariski open dense subset of $Q_b$ for every
$b$ in some nonempty open subset of $S_0$, then
\[
 J=q^{-1}(S_0)\setminus G
\]
is an analytic divisor, possibly empty. Its intersection with
each fibre is proper, so every $G_b$, $b\in S_0$, contains a
Zariski open dense subset of $Q_b$.
\end{proposition}

\begin{proof}
Call a parameter good when its holes have proper complex Zariski
closure. Lemma \ref{holes:bounded-containers} makes the good locus
subanalytic. Choose a finite local stratification of it and its
complement. We first show that no generic real hypersurface wall
can separate an open good region from an open bad region inside
$S_0$.

At a generic point of a proposed wall choose transverse analytic
complex discs, with their real arcs in the wall. Concretely, choose
a real analytic tangent vector field on the wall which is not in
its complex tangent hyperplane, take its local real integral curves,
and complexify those curves. Their imaginary directions are
transverse to the wall, and the resulting family fills a neighborhood.
Lemma
\ref{holes:bounded-containers} gives on each arc a limiting proper
divisor from the good side and its holomorphic continuation on
the disc. All limiting holes on that side are in its support.
The central fibre has a point in $G$, since the wall lies in $S$.
The complement of a proper divisor in a connected smooth
projective variety is connected: a path between two points can
be perturbed off its real-codimension-two strata. Lemma
\ref{holes:tangency-propagation} therefore puts that entire
complement in the wall fibre of $G$.

This also confines holes on both sides near the limiting divisor.
Indeed, any compact subset of the wall fibre outside the divisor
has a neighborhood in $G$; finite coverings and properness make
this uniform on smaller arc intervals. We can now choose
Weierstrass cylinders for the continued divisor. On the real
interval their holes belong to the divisor. Lemma
\ref{holes:interior-trapping}, in the transverse disc, confines
the holes to that divisor on both sides. Varying the transverse
discs through generic wall points yields an open good region on
the proposed bad side, a contradiction.

A connected open real manifold cannot be separated by a
subanalytic subset of real codimension at least two: a compact
path can be perturbed to meet only the hypersurface strata of
a finite stratification. Thus the preceding wall argument shows
that the open good locus is dense in $S_0$. Apply the same
argument to generic hypersurface strata whose two neighboring
open regions are good. It makes their generic parameters good
as well. What remains is a subanalytic subset $E\subset S_0$
of real codimension at least two.

Write $n=\dim_{\mathbb C}Q$ and $f=\dim_{\mathbb C}Q_b$.
Outside $E$ the fibre holes have real dimension at most $2f-2$;
over $E$ they have dimension at most $2f$. The dimension theorem
for subanalytic fibres therefore gives $\dim_{\mathbb R}J\leq2n-2$.
Theorem \ref{holes:thin-analyticity} makes $J$ an analytic divisor.
No fibre is contained in $J$, because every parameter of $S_0$
has a point in $G$. Its intersection with a fibre is consequently
a proper analytic subset of a projective variety, and is algebraic
by Chow's theorem. This proves the final assertion.
\end{proof}

\begin{corollary}[Continuation of a fixed container]
\label{holes:container-propagation}
In Proposition \ref{holes:projective-propagation}, let $D$ be a
relative analytic divisor defined on the full parameter box and
proper in every fibre. If $J\subset D$ over a nonempty open
subregion of $S_0$, then $J\subset D$ over all of $S_0$.
The same conclusion holds for any relative proper analytic
subset in place of $D$.
\end{corollary}

\begin{proof}
Every irreducible component $H$ of $J$ dominates $S_0$. To check
this, $q|_H$ is proper, so its image is analytic. If that image
were proper, dimension of $H$ would force it to be a hypersurface
in $S_0$ with full-dimensional generic fibres in $Q_b$. Those
fibres would equal $Q_b$, contradicting $G_b\ne\varnothing$.
Thus $H$ meets the stated open subregion. The local equations of
$D$ vanish on a nonempty open subset of $H$, and hence on $H$
by the analytic identity theorem. Apply this to every component.
\end{proof}

\begin{corollary}[A section detects the projection domain]
\label{holes:section-test}
Suppose a full-box container $D$ as in Corollary
\ref{holes:container-propagation} is available. A local holomorphic
section $\sigma$ of $q$ avoiding $D$ lies in $G$ at every parameter
of $S_0$. Near a boundary point of $S_0$ that lies in the parameter
box, $S_0$ is locally pseudoconvex wherever such a section exists.
\end{corollary}

\begin{proof}
The first assertion is immediate from containment of the holes.
Since $S=q(G)$ is open, a boundary point of its component $S_0$
cannot belong to $S$: a small connected neighborhood within $S$
would otherwise join it to $S_0$. In submersion coordinates near
$\sigma(b_0)$ choose a connected vertical patch disjoint from
$D$. Over $S_0$ this whole patch is in $G$, while over the
boundary it is disjoint from $G$.

More explicitly, $\sigma^{-1}(G)$ is locally pseudoconvex, since
holomorphic pullback preserves the local Hartogs criterion, and
it is contained in $S$. In a sufficiently small parameter ball,
each component of $S_0$ intersected with that ball is a component
of $\sigma^{-1}(G)$ there: a path in the latter is a path in $S$
and hence cannot leave $S_0$, while every point of $S_0$ belongs
to the pullback. Thus the local pieces of $S_0$ satisfy the same
criterion. This proves the assertion at $b_0$.
\end{proof}
